\section{Introduction}
\label{sec:intro}

Partial observability, game-theoretic reasoning, and long-horizon planning are core challenges in sequential decision-making, yet few existing benchmarks stress all three simultaneously under realistic conditions. Standard testbeds tend to isolate one axis---imperfect-information games emphasize equilibrium computation in short episodes, while open-ended environments test exploration but lack adversarial opponents. Pokémon is an environment that combines all three: competitive battles require reasoning under hidden information against a strategic adversary, while single-player campaigns demand thousands of cumulative decisions spanning exploration, resource management, and combat over extended horizons. With approximately $10^{564}$ possible battle states (see Appendix~\ref{sec:state-space-derivation}), team building across 1,000+ species with diverse movesets and abilities, and a competitive metagame that evolves continuously, Pokémon is more complex and dynamic than most existing benchmarks.

In 2025, Pokémon gained significant interest for evaluating frontier AI systems. Claude Plays Pokémon ~\citep{anthropic2025visible} demonstrated extended thinking capabilities over 35,000 actions to complete a small section of the game, Gemini 2.5 Pro completed the entire game of Pokémon Blue in 406 hours~\citep{gemini2p5report,zhang2025geminiplayspokemon}, and OpenAI's GPT-5 finished the game in 6,470 steps~\citep{engadget2025gptplayspokemon}.
These demonstrations reinforced Pokémon's suitability as an AI testbed, but the efforts were fragmented---different games (Red, Blue, Crystal, Emerald), different harnesses, and different evaluation criteria made meaningful comparison impossible. Was Gemini 3 Pro's 173-hour completion better than Claude Opus 4.6 reaching Victory Road? Did GPT-5's step count account for the same mechanics? By conflating harness with model capability, it became impossible to attribute success to the agent architecture, the underlying model, or hardcoded assumptions that simplified perception. The importance of standardized evaluation in games AI is well established: the Arcade Learning Environment~\citep{bellemare2013arcade} catalyzed a decade of RL progress~\citep{mnih2015human}, while MineRL~\citep{guss2019minerl, milani2023solvingfuzzytaskshuman, shah2022retrospective} demonstrated shared benchmarks for open-ended environments. Pokémon demands---and now receives---similar standardization.

Pokémon also offers a distinctive form of out-of-distribution evaluation. While extensive Pokémon knowledge exists in pretraining corpora---move types, damage formulas, competitive tier lists---translating that latent knowledge into effective multi-turn sequential decision-making under partial observability is fundamentally different from the recognition and retrieval tasks where data contamination typically inflates performance. Moreover, the competitive metagame shifts continuously as the player community develops new strategies and the game's governing body rebalances tiers---creating natural distribution shifts that test an agent's ability to adapt rather than memorize.

We present the \pokeagent Challenge, a standardized evaluation framework for Pokémon-playing AI agents. The benchmark features two complementary tracks: \textbf{Competitive Battling} evaluates strategic reasoning under partial observability in two-player competitive Pokémon, while the \textbf{RPG Speedrunning} tests long-horizon planning in completing Pokémon Emerald as quickly as possible.

The NeurIPS 2025 \pokeagent Challenge confirmed the benchmark's difficulty and drew strong community engagement: 100+ teams submitted solutions across both tracks and 650+ researchers joined the competition Discord for technical exchange. The competition produced novel methods---including Scripted Policy Distillation for RPG play and iterative offline RL with dynamic data weighting---and served as the first large-scale competitive testbed for approaches such as root-parallelized MCTS in imperfect-information battling~\citep{FoulPlay}, while revealing a capability hierarchy: specialist RL and search methods dominated LLM approaches in Competitive Battling, and no raw frontier model achieved non-trivial progress in Speedrunning without a sophisticated harness. These gaps remain far from closed.

Our contributions include: (1)~the \pokeagent Challenge, a two-track evaluation framework pairing competitive battling (via Pokémon Showdown) with RPG speedrunning (via Pokémon Emerald), with standardized infrastructure for fair comparison across RL, LLM, and hybrid approaches; (2) the largest publicly available Pokémon battle dataset---comprising 4M human demonstrations and 18M synthetic battles, plus 200K+ curated competitive teams; (3)~baselines spanning heuristic bots, RL agents, and harness LLM agents, alongside the first open-source multi-agent orchestration system for long-horizon RPG play; (4)~empirical validation through the NeurIPS 2025 \pokeagent Challenge (100+ competitors, 100K+ battles, top methods in Appendix~\ref{sec:participant-methodologies}), with results revealing substantial gaps between generalist LLM, specialist RL, and elite human performance, and orthogonality analysis showing that Pokémon battling captures capabilities not predicted by the 49 benchmarks in the BenchPress evaluation matrix~\citep{papailiopoulos2026benchpress}; and (5)~a living benchmark with a live Battling and Speedrunning leaderboard and self-contained Track~2 evaluation at \url{https://pokeagentchallenge.com}.